\documentclass[prep,both]{debate}
\motion{更应该用现实 / 更应该用故事向孩子解释死亡}
\meta{赛制}{中国科学技术大学新生辩论赛（默认赛制，四对四）}
\meta{生成日期}{2026-09-15}
\begin{document}
\debatetitle
\begin{summary}
\begin{fields}
\claim[持方优劣评估]{\textbf{反方（故事）较劣势}，劣势集中在「论据可得性」与「论证义务」两个维度。发展心理学的实证证据与所有权威机构的家长指南几乎一边倒地站在「直接、诚实、生物学」一侧，反方拿不到同等级的正面证据。反方唯一的天然优势是评委直觉：多数人自己是听着「变成星星了」长大的。反方的打法是把战场收窄到「孩子几岁」与「解释的目的是什么」，抢先把「故事」从「哄骗」里摘出来，并在价值层完成反超}
\thesis{正方}{孩子来问死亡的时候问的是事实，只有先把事实说清楚，后面的温柔才有地方落脚 }
\thesis{反方}{解释成功与否不看大人说了什么，而看孩子接住了什么，而在孩子最需要解释的那个年纪，他能接住的是故事 }
\end{fields}
\end{summary}
\section{一、赛制要点}
\begin{flowtable}
\block{A}
\flow{1 正一立论 / 2 反四质询正一 / 3 反一立论 / 4 正四质询反一}{180+120+180+120 秒}{正一、反四、反一、正四}{立论清楚、判准站得住；质询拿到对方在「年龄」或「场景」上的一个明确承认}
\block{B}
\flow{6 反二驳论 / 7 正二驳论 / 8 二辩对辩}{120+120+各 90 秒}{正二、反二}{把对方立论的核心问题讲给评委听；对辩不丢判准}
\block{C}
\flow{10 正三盘问 / 11 反三盘问 / 12 正三小结 / 13 反三小结}{各 90 秒}{正三、反三、双方二辩四辩}{盘问链完成；被盘问时二辩与四辩口径不打架}
\block{D}
\flow{15 自由辩论}{双方各 240 秒}{全员}{主战场推进到位；对方必答问题都正面回应过}
\block{E}
\flow{17 反四结辩 / 18 正四结辩}{各 210 秒}{反四、正四}{回到判准与一句话立论，处理掉对方全场最强的一个点}
\end{flowtable}
计时方式对策略的影响：质询是双边计时，被质询方拖时间就是在吃质询方的表，所以这道题里两位一辩被问到「几岁」「哪个场景」时，回答必须短而且第一句就给出立场，不要展开。盘问是单边计时，对方回答一般不计时，所以三辩可以放对方多说几句，让二辩与四辩的口径差异自己露出来。自由辩双方各 240 秒，这道题的交锋点密集，发言一定要一次一个点，说完就坐。

\section{二、辩题性质与争议焦点}
\begin{fields}
\claim[辩题性质]{\textbf{比较题}，带价值倡导色彩。“更应该”意味着双方都承认对方的方法有价值，争的是在同一个解释场合里哪一种\textbf{作为主导方式、作为默认起点}。任何一方试图证明“对方那一套完全没用”，都是在打自己不需要打的仗}
\field{正方倾向把题目解释成}{一道关于“解释的准确性与后果”的事实兼价值题。正方希望评委问的是“哪种方式让孩子真的知道发生了什么，并且不因此受伤”
}
\field{反方倾向把题目解释成}{一道关于“沟通对象的接收能力”的题。反方希望评委问的是“在孩子那个年纪，哪种方式的解释真的进得去”
}
\field{争议焦点}{
\begin{fields}
\field{年龄之争}{把默认的“孩子”设在几岁。设在三到五岁对反方有利，设在六到十二岁对正方有利
}
\field{场景之争}{默认场景是“孩子某天忽然问起”（预先的教育），还是“亲人今天刚去世”。前者对正方有利，后者对反方有利
}
\field{准确性是否被占用}{讲故事会不会挤掉孩子对死亡的正确理解。这一条上双方各有一篇顶刊研究，是全场最硬的一次对撞（见第五节论点二与第六节论点二）
}
\end{fields}
}
\end{fields}
\section{三、关键词定义}
\subsection{3.1 现实}
\begin{fields}
\field{调研}{《现代汉语词典》（第 7 版）“现实”条：客观存在的事物；合于客观情况。在本题语境下，“用现实解释”指以死亡的客观属性作为解释内容。医学与发展心理学界通用的死亡概念成分，来自 Speece 与 Brent（1984，《Child Development》55 卷 5 期 1671–1686 页）：不可逆性、丧失功能性、普遍性；他们 1992 年又补入因果性。英国国民保健署（NHS）“Children and bereavement”页（最后审阅 2023 年 1 月 9 日）的表述是：直接、诚实、开放的沟通比隐瞒真相更有帮助
}
\begin{procard}{正方有利定义}\definition{现实指“与事实一致、可核对的解释”，不限定语气。}这样一来“温柔地说清楚”也算现实，正方就同时占住了准确与温柔两块地。\sub{反方如何反驳}{这个定义把“怎么说”整个抽掉了，只剩“说什么”，等于把辩题改成“要不要说真话”，反方持方在这个定义下不可能成立}\end{procard}
\begin{concard}{反方有利定义}\definition{现实指“以生物学事实的陈述为内容的解释”，即告诉孩子器官停止工作、身体不再运转、人不会回来。}\sub{正方如何反驳}{这是把正方窄化成“医学科普”，辩题里没有任何词要求解释只能是名词解释；用这个定义等于替正方选了一个最难守的版本}\end{concard}
\begin{neutralcard}{中立定义}\textbf{用现实解释死亡，指以死亡的事实属性为解释的主要内容，核心是三句话：人死了不会再回来，身体停止运作，每个人最后都会死。语气可以温柔，也可以严肃，但字面内容与事实一致。}\sub{合理性}{贴近常识（普通人说“跟孩子讲实话”就是这个意思），也贴近权威机构的家长指南所指的那件事，并且把语气问题留给双方去争，辩论才有得打}\end{neutralcard}
\end{fields}
\subsection{3.2 故事}
\begin{fields}
\field{调研}{《现代汉语词典》（第 7 版）“故事”条：真实的或虚构的有人物有情节的事情。在儿童死亡教育的实务里，“故事”指绘本叙事、比喻与拟人化说法，例如“变成了星星”“去了另一个地方”，或《獾的礼物》《爷爷变成了幽灵》这类文本。美国儿科学会（HealthyChildren.org，“How Children Understand Death”，Schonfeld 与 Nasir 撰，2024 年 6 月 26 日更新）把“睡着了”“去了很远的地方”列为应避免的含糊说法
}
\begin{procard}{正方有利定义}\definition{故事指“用与事实不符的说法代替事实的解释”。}\sub{反方如何反驳}{这是标准的定义杀，把“故事”等同于“欺骗”，反方持方在这个定义下不可能成立，评委会判正方在逃避辩论。\textbf{正方绝对不要用这一版}}\end{procard}
\begin{concard}{反方有利定义}\definition{故事指“以叙事形式承载死亡的意义，可以完全不违背事实”。}\sub{正方如何反驳}{那你已经在用现实了，只是加了包装；如果故事的字面内容必须与事实一致，双方的分歧就只剩修辞，辩题失去意义}\end{concard}
\begin{neutralcard}{中立定义}\textbf{用故事解释死亡，指以叙事、比喻、拟人这类非直陈形式作为解释的主要载体，其字面内容不要求与事实一一对应。}\sub{合理性}{这一版既容纳了“变成星星”这类明显偏离事实的说法，也容纳了《獾的礼物》这类不违背事实的叙事，双方都有空间；同时它不预设故事就是谎言}\end{neutralcard}
\end{fields}
\subsection{3.3 孩子}
\begin{fields}
\field{调研}{发展心理学的关键节点在这里。Speece 与 Brent（1984）综述四十余项研究后得出：在现代城市工业社会，多数健康儿童在\textbf{五到七岁之间}获得对不可逆性、丧失功能性、普遍性的理解。Menendez、Hernandez 与 Rosengren（2020，《Child Development Perspectives》，DOI 10.1111/cdep.12357）的综述把节点写得更细：五岁前开始理解普遍性，随后是终结性；五岁时多数孩子理解身体机能停止；六岁时理解死亡可以由多种原因造成。Renaud 等（2015）在加拿大的网络调查里，家长报告孩子最早在\textbf{三到三岁半}就开始提问
}
\begin{procard}{正方有利定义}\definition{孩子指“已经会开口提问的儿童”，重心放在六岁以上。}\sub{有利之处}{六岁之后生物学理解已经基本成形，正方的所有证据都在这个区间最强。}\sub{反方如何反驳}{孩子第一次问死亡是三岁半，你把最需要解释的那一段切掉了}\end{procard}
\begin{concard}{反方有利定义}\definition{孩子指“学龄前儿童”，重心放在三到五岁。}\sub{有利之处}{这个年纪的生物学理解还没成形，反方的“接收不了”机制最成立。}\sub{正方如何反驳}{辩题没说学龄前；而且真正频繁遇到丧亲、需要正式解释的年龄段比这更宽}\end{concard}
\begin{neutralcard}{中立定义}\textbf{孩子指三到十二岁的儿童，讨论重心放在三到八岁}，理由是提问最早出现在三岁半，而死亡概念在五到七岁之间成形，这一段正是“解释”这件事真正发生的时候。\textbf{双方各准备一段三十秒的年龄口径}（见第十节口径表），被问到就照着说，不要现场找\end{neutralcard}
\end{fields}
\subsection{3.4 隐含要素}
\begin{fields}
\field{主体}{做解释的是家长、老师或其他成年抚养者；受影响的是三到十二岁儿童
}
\field{范围}{当代中国的普通家庭为主，可以引用跨文化研究做机制说明，但要交代样本来自哪里
}
\field{时间尺度}{短期看孩子当下是否被安抚，长期看孩子最终形成的死亡概念与信任关系。两个尺度都要算
}
\field{比较对象}{现实与故事二者相比，不是“现实”与“什么都不说”相比。任何一方把对手换成“回避不谈”（“你还小，长大就懂了”），都是在打稻草人
}
\field{场景}{两类，预先的教育（孩子某天忽然问起）与丧亲当日（亲人刚刚去世）。\textbf{双方都必须能处理两类}，只是各自会把不同的一类设为默认
}
\end{fields}
\section{四、判准与论证义务}
\begin{fields}
\claim[判准是什么]{这道题争的是一种解释方法的优先性，所以判准要同时装下“说得对不对”和“孩子受不受得住”，并且要规定两者冲突时怎么办}
\begin{procard}{正方有利判准}\definition{以“孩子最终获得的死亡概念是否正确”为准。}\sub{有利之处}{所有实证证据都在这条线上，正方几乎自动获胜。}\sub{反方如何反驳}{这条标准把解释的另一半目的整个删掉了。一个孩子在葬礼当天需要的不只是正确的知识；按这条标准，最好的做法是给四岁孩子放一段解剖学教学视频，评委自己都不会同意}\end{procard}
\begin{concard}{反方有利判准}\definition{以“孩子当下是否被安抚、是否愿意继续跟大人谈”为准。}\sub{有利之处}{故事在情绪安抚上确实更顺。}\sub{正方如何反驳}{这条标准把“说的是不是真的”排除在外，那么任何有效的善意谎言都能过关，包括“爷爷出差了，过几年就回来”。评委不会接受一条让欺骗自动得分的标准}\end{concard}
\begin{neutralcard}{中立判准}\textbf{看哪一种方式更能让孩子在这次解释之后，既正确地知道发生了什么，又承受得住它、并且愿意继续跟大人谈这件事。} 两个指标：认知上的准确，情感上的可承受。两者冲突时，看\textbf{哪一方的失误更难挽回}。\sub{合理性}{它照顾了题目的主体（孩子不是知识容器，也不是只需要哄的对象），符合比较题的形态（两项指标上做净比较），并且真的能区分双方}\end{neutralcard}
\begin{fields}
\field{正方需要证明}{以现实为主导的解释，在这两个指标上整体更优；并且要正面处理“孩子承受不住”这个担心，说明它被高估在哪里
}
\field{反方需要证明}{以故事为主导的解释，在这两个指标上整体更优；并且要正面处理“故事会让孩子误解事实”这个担心，说明这种误解要么不发生，要么可修正、代价小
}
\end{fields}
\end{fields}
\prosection{五、正方论点}
\prosubsection{论点一：孩子问的是事实}
\begin{fields}
\claim[主张]{孩子来问死亡的时候，问的是事实问题；用故事回答，孩子拿到的不是答案，是一个他会按字面理解的新说法。}
\begin{mechanism}[机制]
\step 孩子问“为什么会死”“爷爷还会回来吗”
\step 大人答“去了天上”“睡着了”
\step 孩子按字面接收
\step 产生新的恐惧或错误预期（怕睡觉、怕出门、等人回来）。

\end{mechanism}
\begin{evidence}{学理}
Speece 与 Brent（1984，1992）的死亡概念成分论。\sub{核心主张}{死亡概念不是一个整块，而是不可逆性、丧失功能性、普遍性、因果性四个可以分别获得的成分。}\sub{支撑}{它解释了机制的第一步，孩子的提问正是冲着这些成分来的，尤其是最晚形成的因果性。}\sub{边界}{这套框架描述的是概念如何获得，不直接推出“该怎么说”，正方不要把它当成结论。}\sub{出处}{Speece M W, Brent S B (1984)，《Children's Concept of Death: A Review of Three Components of a Death Concept》，《Child Development》55(5): 1671–1686。}
\end{evidence}
\begin{evidence}{数据}[已核实]
Gutiérrez 等（2020）对墨西哥普埃布拉 53 个家庭、61 名三岁半到六岁十一个月儿童的混合方法研究：孩子关于死亡的一般提问中 61.2\% 是关于死因的，\textbf{没有一个孩子的提问提到宗教}；而家长的回答里有 22.0\% 含宗教内容。\sub{来源}{Gutiérrez I T, Menendez D, Jiang M J, Hernandez I G, Miller P, Rosengren K S (2020)，《Embracing Death: Mexican Parent and Child Perspectives on Death》，《Child Development》，DOI 10.1111/cdev.13263。}\sub{方法}{53 个家庭、61 名儿童（平均 5.1 岁）的结构化访谈加家长问卷，量化与民族志结合。}\sub{局限}{样本在墨西哥，死亡在当地文化中的地位与中国不同；它证明的是“错配存在”，不是“错配有害”。}\sub{状态说明}{。}
\end{evidence}
\begin{evidence}{案例}
美国儿科学会在其家长指南里直接点名“爷爷睡着了”“去了一趟远门”这类说法，理由是孩子按字面理解，可能因此害怕睡觉、害怕出门；儿童丧亲慈善机构 Winston's Wish 举的例子更直白：如果说“我们把妈妈弄丢了”，孩子会问“那我们为什么不去找她”。\sub{代表性}{这是实务界反复遇到的典型情况，不是极端个例。}\sub{出处}{HealthyChildren.org（美国儿科学会），《How Children Understand Death: What to Say When a Loved One Dies》，Schonfeld D J 与 Nasir A 撰，2024 年 6 月 26 日更新；Winston's Wish，《Telling a child someone has died》。}
\end{evidence}
\closing[回扣判准]{认知指标上孩子没拿到答案，情感指标上还多了一份新的害怕，两项都是负的。}
\begin{clash}[对方最强攻击 → 我方回应]
反方会说“你举的是糟糕的故事，不是故事本身”。
\end{clash}
\end{fields}
\prosubsection{论点二：现实是地基，故事是盖在上面的楼}
\begin{fields}
\claim[主张]{孩子是先有生物学理解，再用它去约束比喻和来世的说法，不是反过来；而且越早获得这份理解，越不怕死。}
\begin{mechanism}[机制]
\step 懂得身体靠什么维持运转
\step 懂得死亡意味着什么
\step 恐惧下降，并且有了校对其他说法的尺子。

\end{mechanism}
\begin{evidence}{学理}
Slaughter 的生机论学习框架，以及 Menendez 等（2020）综述给出的顺序结论：儿童先形成生物学理解，再用它约束宗教理解。\sub{核心主张}{孩子对生命的因果理解是死亡概念的支架。}\sub{支撑}{它解释了机制的第一步与第二步。}\sub{边界}{这个顺序是描述性的，不能直接推出“越早讲越好”，正方要靠训练实验补上因果。}\sub{出处}{Menendez D, Hernandez I G, Rosengren K S (2020)，《Children's Emerging Understanding of Death》，《Child Development Perspectives》，DOI 10.1111/cdep.12357。}
\end{evidence}
\begin{evidence}{数据}[已核实]
Slaughter 与 Lyons（2003）的\textbf{训练实验}：60 名三岁七个月到五岁十一个月的学龄前儿童，其中 40 人接受关于人体如何运转以维持生命的教学，20 人为对照组；训练后，训练组在人体功能与死亡两方面的理解都出现显著发展。\sub{来源}{Slaughter V, Lyons M (2003)，《Learning about life and death in early childhood》，《Cognitive Psychology》46(1): 1–30，DOI 10.1016/s0010-0285(02)00504-2。}\sub{方法}{前测、分组教学、后测的准实验设计，这是因果证据不是相关证据。}\sub{局限}{样本 60 人；教的是身体功能而不是丧亲场景的对话，推广到丧亲当日要谨慎。}\sub{状态说明}{。}
\end{evidence}
\begin{evidence}{数据}[已核实]
Slaughter 与 Griffiths（2007）对 90 名四到八岁儿童的访谈研究：在控制年龄与一般焦虑之后，越成熟的生物学死亡理解对应越低的死亡恐惧。\sub{来源}{Slaughter V, Griffiths M (2007)，《Death understanding and fear of death in young children》，《Clinical Child Psychology and Psychiatry》12(4): 525–535。}\sub{方法}{横断面访谈加家长报告的一般焦虑量表，回归分析。}\sub{局限}{相关不是因果，可能是不那么焦虑的孩子更愿意正面回答死亡问题。}\sub{状态说明}{。}
\end{evidence}
\begin{evidence}{数据}[已核实]
Ellis、Dowrick 与 Lloyd-Williams（2013）对 33 位十八岁前失去父母的成年人的叙事研究发现，没有在合适的时点、按孩子的理解水平给出清楚诚实的信息，被受访者认为对其成年后的信任感、自我价值、孤独感有负面影响。\sub{来源}{《The long-term impact of early parental death: lessons from a narrative study》，《Journal of the Royal Society of Medicine》106(2): 57–67。}\sub{方法}{33 人（7 男 26 女）的书面与口述叙事，质性叙事分析。}\sub{局限}{回溯性自述，样本小且来自英格兰西北部。}\sub{状态说明}{。}
\end{evidence}
\begin{evidence}{案例}
《一片叶子落下来》把死亡讲成“另一种形式的新生”。这句话在成年人耳朵里是哲学，在一个四岁孩子耳朵里是一条关于世界如何运作的信息，而且它和不可逆性正好相反。\sub{出处}{Buscaglia L（1982），《The Fall of Freddie the Leaf》，中译本《一片叶子落下来》，任溶溶译。}
\end{evidence}
\closing[回扣判准]{认知指标与情感指标同向为正，这是正方最关键的一句话。}
\begin{clash}[对方最强攻击 → 我方回应]
反方会用 Panagiotaki 等（2018）指出成人组已经没有共存式思维，说明听故事长大的孩子最终照样懂。
\end{clash}
\end{fields}
\consection{六、反方论点}
\consubsection{论点一：在那个年纪，故事是唯一进得去的通道}
\begin{fields}
\claim[主张]{解释算不算成功，不看大人说出了什么，看孩子接住了什么；在孩子最需要解释的三到六岁，能被接住的形式是叙事。}
\begin{mechanism}[机制]
\step 这个年纪的孩子以叙事方式理解世界
\step 故事把无法想象的事变成有人物、有时间、有结局的事
\step 孩子能复述、能记住、能带着走。

\end{mechanism}
\begin{evidence}{学理}
Bruner（1986）的两种思维模式：范式模式用范畴、逻辑与可检验命题解释世界，叙事模式用有意图的人物、麻烦与结局解释世界；两者互不可还原。\sub{核心主张}{叙事不是范式思维的简化版，它是另一套完整的理解方式。}\sub{支撑}{它解释了机制的第一步，说明为什么“讲得更准确”不等于“讲得更进得去”。}\sub{边界}{这是认识论框架不是实证研究，反方只能用它解释机制，不能拿它当证据。}\sub{出处}{Bruner J（1986），《Actual Minds, Possible Worlds》，Harvard University Press。}
\end{evidence}
\begin{evidence}{数据}[已核实]
Liu 与 Liu（2024）对河南省一所公立幼儿园 114 名五到六岁儿童的研究显示，死亡概念五个子成分的成熟率分别是：不可逆性 67.8\%、因果性 60.7\%、非功能性 57.1\%、普遍性 55.3\%、死亡判定 53.6\%。也就是说，在孩子最常追问死亡的年纪，即使最容易的那个成分，也有三分之一的孩子还没掌握。\sub{来源}{Liu J, Liu F (2024)，《Research on the development of the concept of death in children and its influencing factors》，《Frontiers in Psychology》15，DOI 10.3389/fpsyg.2024.1376253。}\sub{方法}{114 名儿童（59 男 55 女）的结构化访谈加母亲问卷，偶遇抽样，访谈有效率 97.56\%。}\sub{局限}{单园单地区，偶遇抽样，不能代表全国。}\sub{状态说明}{。}
\end{evidence}
\begin{evidence}{数据}[已核实]
Mar、Li、Nguyen 与 Ta（2021）的元分析，基于 75 个以上独立样本、33000 多名参与者：故事比说明文更容易被理解，也被记得更牢，结果稳健，不受各类研究特征调节。\sub{来源}{《Memory and comprehension of narrative versus expository texts: A meta-analysis》，《Psychonomic Bulletin \& Review》28(3): 732–749。}\sub{方法}{对比叙事文本与说明文本在记忆与理解上差异的实验的元分析。}\sub{局限}{\textbf{测的是文本的理解与记忆，不是死亡观的形成，参与者也以较大年龄为主}。反方只能用它支撑“叙事作为传递形式更有效”这一步机制，不能用它证明“故事讲的内容更对”。稿件中必须自己先说出这条局限。}\sub{状态说明}{。}
\end{evidence}
\begin{evidence}{案例}
《獾的礼物》。獾老了，它死了，书里没有让它复活，也没说它去了别处。朋友们靠着回忆它教过的每一件事活下去。孩子从这本书里带走的不是“心脏停止工作”，是“獾教会鼹鼠剪一串纸手拉手的小人”。这说明一个不违背任何事实的故事，照样能把死亡讲完。\sub{出处}{Varley S（1984），《Badger's Parting Gifts》，Andersen Press，1985 年获 Mother Goose Award。}
\end{evidence}
\closing[回扣判准]{认知指标上孩子拿到了他能复述的东西，情感指标上他拿到了一个能承受的形式。}
\begin{clash}[对方最强攻击 → 我方回应]
正方会举 Ganea 等（2014）：拟人化绘本让孩子学到的新事实更少，还会把人的特征安到真动物身上。
\end{clash}
\end{fields}
\consubsection{论点二：故事不占用准确性，却多给了一样东西}
\begin{fields}
\claim[主张]{孩子完全可以一边知道人死了不会回来，一边在故事里保留这段关系；而后面这件事，只有故事做得到。}
\begin{mechanism}[机制]
\step 生物学理解与叙事性理解在孩子心里可以并存
\step 故事不占用准确性的份额
\step 故事额外提供意义与联结
\step 哀伤有地方安放。

\end{mechanism}
\begin{evidence}{学理}
Klass、Silverman 与 Nickman（1996）提出的持续联结理论。\sub{核心主张}{与逝者维持联结不是病理，而是正常且具有保护作用的哀伤方式，它取代了此前“哀伤的功能是切断联结”的旧模型。}\sub{支撑}{它解释了机制的第三步，说明“多给的那一样东西”到底是什么。}\sub{边界}{持续联结讲的是哀伤过程，不直接推出“应该用比喻对孩子说话”，反方要靠数据补上那一步。}\sub{出处}{Klass D, Silverman P R, Nickman S（1996），《Continuing Bonds: New Understandings of Grief》。}
\end{evidence}
\begin{evidence}{数据}[已核实]
Gutiérrez 等（2020）的同一份研究：81.9\% 的孩子相信亡故的亲人在亡灵节回来探望，78.6\% 相信亡者吃掉了供品；与此同时，\textbf{全部孩子都表现出至少部分生物学死亡理解}，并且父母自报的宗教与灵性思考，对普遍性、终结性、非功能性、因果性四个子成分\textbf{均无显著影响}。来源同第五节论点一。方法同上。\sub{局限}{墨西哥样本，死亡在当地文化中的地位特殊；且这是横断面数据，不能证明故事带来了好处，只能证明它没有挤掉准确性。}\sub{状态说明}{。}
\end{evidence}
\begin{evidence}{案例}
中国的清明。一家人在墓前摆上三只碗，说一句“爷爷你吃”。没有哪个大人因此真的以为爷爷会回来吃，也没有哪个孩子因此在科学课上答错人死了会不会复活。这就是共存：一套用来知道，一套用来相处。
\end{evidence}
\closing[回扣判准]{认知指标没有被牺牲，情感指标上多了一项，净差为正。}
\begin{clash}[对方最强攻击 → 我方回应]
正方会用 Panagiotaki 等（2018）：六到七岁与八到九岁两组里，非宗教家庭的孩子在“普遍性”上更常给出生物学正确答案（均值 1.67 比 0.83，以及 1.65 比 0.89）；八到十一岁中宗教家庭孩子的生物学模型只占 23\%，非宗教家庭占 77\%。回应两层。第一，那项研究测的是\textbf{父母信不信来世}，不是\textbf{父母讲不讲故事}，这两件事不能划等号，一个无神论家长照样会说“姥姥变成星星了”。第二，同一篇论文的成人组\textbf{没有出现共存式思维}，只剩下生物学解释与超自然解释的清楚区分。也就是说，这笔账后来是被还清的；而故事提供的那份联结不会被收回。
\end{clash}
\end{fields}
\consubsection{被淘汰的候选论点（进反驳库，对手可能会用）}
正方淘汰：故事被拆穿以后成本是信任；孩子反正会从媒体接触死亡（Bridgewater 等 2019 报告 75\% 的儿童动画片里有死亡），不如由大人先讲对；中国家长的回避本身造成了死亡教育的空白。反方淘汰：丧亲当天讲生物学是二次伤害；大人自己也过清明、写墓志铭，只要求孩子接受现实是双标；绘本疗法在儿童哀伤辅导中的应用。

\section{七、正方反驳库（针对反方全部可能论点）}
\begin{rebuttals}
\rebuttal{1}{三到六岁的孩子接不住生物学解释，故事是唯一通道}{接不住完整的生物学，不等于接不住“不会回来”这五个字}{打机制。Slaughter 与 Lyons 的训练实验里，三岁七个月起的孩子被教了人体如何维持生命之后，对死亡的理解显著提升。所以“这个年纪听不懂”是被高估的。何况我方要求的不是讲循环系统，是讲清楚不可逆性}{追问“你有没有证据说四岁孩子能懂不可逆性” → 回应：Menendez 等 2020 的综述写明，五岁前孩子就开始理解普遍性与终结性。而且对方自己引的河南研究里，五到六岁孩子不可逆性成熟率 67.8\%，是五项里最高的，正好说明这一项最先能接住}{最终}
\rebuttal{2}{故事与生物学理解可以共存，不占用准确性}{共存不是免费的，你自己的数据算的是短期，代价在后面}{打论据。Panagiotaki 等 2018 的英国样本里，八到十一岁中宗教家庭孩子的生物学模型只有 23\%，非宗教家庭是 77\%；六到七岁与八到九岁两组在普遍性上也都是非宗教家庭更准确。对方那份墨西哥研究的样本上限只到六岁十一个月，年龄没跨到差异显现的区间}{追问“那测的是父母信仰不是讲故事” → 回应：讲不讲来世的故事，与信不信来世高度重合；退一步说，就算只是相关，举证责任在主张“无害”的一方}{最终}
\rebuttal{3}{叙事比说明文更好懂更好记}{这份元分析测的是读课文，不是懂死亡}{打论据的适用范围。Mar 等 2021 测的是文本理解与回忆，参与者以较大年龄为主。把它推到“四岁孩子怎么理解爷爷去世”，中间跨了两步。而且好记不等于记对，Ganea 等 2014 恰恰显示拟人化叙事让孩子学到的事实更少}{追问“你是说叙事没有优势吗” → 回应：我方承认叙事在传递上更有效。正因为有效，讲错的东西也会被更牢地记住}{最终}
\rebuttal{4}{故事维系了与逝者的联结，这对哀伤有保护作用}{联结靠的是记忆，不靠比喻}{打机制。持续联结理论说的是保留关系，没有说保留关系必须靠一个不真实的字面说法。《獾的礼物》里的联结是靠回忆建立的，它一句谎话都没说}{追问“那你承认《獾的礼物》是好方法” → 回应：承认，而且它正是我方说的，那本书从头到尾没有否认獾死了。对方要论证的是“字面不必与事实对应”这件事更好，不是“叙事形式”更好}{最终}
\rebuttal{5}{丧亲当天对孩子讲生物学是二次伤害}{我方从没主张在葬礼上讲解剖学}{打定义。中立定义里的现实是三句话：不会再回来、身体停止运作、每个人都会死，语气可以温柔。Ellis 等 2013 对 33 位童年丧亲者的叙事研究发现，没有在合适时点给出清楚诚实的信息，反而对成年后的信任与自我价值有负面影响}{追问“那孩子当场崩溃怎么办” → 回应：抱着他。安慰不需要谎言做前提，这一点对方一直在混为一谈}{淘汰}
\rebuttal{6}{大人自己也过清明、也说“在天上看着你”，只要求孩子接受现实是双标}{大人知道那是仪式，孩子不知道}{打机制。仪式对成年人是一种表达方式，对一个还没掌握不可逆性的孩子却是一条字面信息。Panagiotaki 等 2018 里有受访成人直接承认，自己并不相信来世，却仍然告诉孩子“奶奶在天上还看得见你”，因为这样“威胁更小、更安慰人”}{追问“那你反对清明吗” → 回应：不反对。我方反对的是把仪式当成解释。孩子先要知道爷爷不会回来，再去墓前摆那只碗，两件事的顺序不能颠倒}{常见}
\rebuttal{7}{故事让孩子愿意继续谈，直陈会让对话终结}{让对话终结的是回避，不是诚实}{打判准。真正终结对话的是“你还小，长大就懂了”。NHS 的家长指南写得很清楚：直接、诚实、开放的沟通比隐瞒真相更有帮助}{追问“孩子被吓哭了还怎么谈” → 回应：哭不是对话终结，是对话开始。真正的终结是孩子发现大人说的不算数}{常见}
\rebuttal{8}{绘本在儿童哀伤辅导里被广泛使用，说明故事有效}{广泛使用不等于有效，这一块的证据很薄}{打论据。绘本疗法针对丧亲儿童的随机对照证据数量少、异质性高，目前无法给出确定性结论。相比之下我方有训练实验。而且被辅导师推荐的绘本，绝大多数恰恰是不违背事实的那一类}{追问“你是说绘本没用吗” → 回应：不是。我方说的是它作为辅助有用，作为主导缺乏证据}{常见}
\end{rebuttals}
\section{八、反方反驳库（针对正方全部可能论点）}
\begin{rebuttals}
\rebuttal{1}{孩子问的是事实，用故事回答是答非所问}{你自己引的数据说的是父母“补充”了宗教，不是“替换”了事实}{打论据。Gutiérrez 等 2020 的原文里写明，许多父母是在回答孩子问题的基础上\textbf{扩展}到来世与仪式。孩子问死因，父母答了死因（55.0\%），另外加了宗教（22.0\%）。那不是答非所问，那是答完了还多说了一句}{追问“多说的那句会被孩子当真” → 回应：同一份研究里，全部孩子仍然表现出生物学死亡理解，父母的宗教与灵性思考对四个子成分均无显著影响}{最终}
\rebuttal{2}{孩子会按字面理解比喻，产生新的恐惧}{那是差劲的比喻的问题，不是故事的问题}{打定义。对方全场举的例子只有“睡着了”“弄丢了”这两个，这两个连故事都算不上，是敷衍。请对方用《獾的礼物》举一次例。做不到，说明对方反对的其实是敷衍，不是叙事}{追问“你怎么保证家长讲的是好故事” → 回应：同理，你也无法保证家长讲的现实不是“人死了就是没了，别问了”。方法之争比的是方法本身能达到的上限}{最终}
\rebuttal{3}{先有生物学理解，才能约束比喻；恐惧也会下降}{你的训练实验教的是身体，不是解释死亡}{打论据的适用范围。Slaughter 与 Lyons 2003 教的是人体如何运转，样本 60 人，场景是课堂不是丧亲。把它当成“应该用现实向孩子解释死亡”的直接证据，中间跨了一大步}{追问“那你承不承认懂生物学的孩子更不怕死” → 回应：承认相关，那是 Slaughter 与 Griffiths 2007 的横断面数据，相关不是因果。也可能是本来就不那么焦虑的孩子更愿意正面谈死亡}{最终}
\rebuttal{4}{故事被拆穿以后，孩子会不再信任大人}{孩子并不把故事当成可被拆穿的事实陈述}{打机制。墨西哥的孩子既相信亡者回来探望，又知道死亡不可逆，两套并存而没有崩塌。信任危机来自“被隐瞒了发生的事”，不来自“听过一个故事”}{追问“那孩子长大发现没有天堂呢” → 回应：Panagiotaki 等 2018 的成人组没有出现共存式思维，只有清楚的两套解释。他们没有崩溃，他们长大了}{淘汰}
\rebuttal{5}{权威机构的家长指南都站在诚实一侧}{那些指南反对的是含糊，不是叙事}{打论据。美国儿科学会点名的是“睡着了”“去了很远的地方”，英国那几家慈善机构点名的是“弄丢了”“走了”。它们要求的是别用含糊的替代词，同时这些机构自己就在推荐绘本。Child Bereavement UK 甚至为《獾的礼物》配了家长阅读指南}{追问“那它们也说要用‘死了’这个词” → 回应：我方不反对说“死了”。我方主张的是，说完这两个字之后，靠什么把这件事讲下去}{最终}
\rebuttal{6}{中国家长回避死亡教育，所以更该补现实这一课}{回避的反面是开口，不是“只讲生物学”}{打定义。对方把“故事”和“回避”混成了一件事。真正的回避是“你还小，长大就懂了”。一个愿意坐下来给孩子读《爷爷变成了幽灵》的家长，恰恰是没有回避的那个}{追问“绘本里说爷爷变成幽灵回来了，这不是回避是什么” → 回应：那本书的结尾是爷爷说完再见就走了，再也没回来。它讲的正是告别}{常见}
\rebuttal{7}{五到七岁孩子已经能理解，不必绕弯}{你引的每一份数据都说“多数”，剩下那部分孩子怎么办}{打论据。Liu 与 Liu 2024 的河南样本里，五到六岁孩子成熟率最高的不可逆性也只有 67.8\%，最低的死亡判定只有 53.6\%。面对一个具体的孩子，大人不知道他在哪一边。故事的好处是，接不住的孩子也能接住一半}{追问“接住一半的是错的一半怎么办” → 回应：《獾的礼物》里没有错的一半。这个问题只对偷懒的比喻成立}{最终}
\rebuttal{8}{解释的目的就是让孩子知道真相}{那是判准之争，而且这条判准评委不会接受}{打判准。按这条标准，最好的做法是给四岁孩子放一段解剖视频。孩子来问死亡，从来不只是在要一个知识，他是在问“我还会不会被爱”“你会不会也走”。这两个问题，事实回答不了}{追问“你是说真相不重要” → 回应：重要，所以我方从不主张说假话。我方主张的是，真相要有一个孩子背得动的形状}{常见}
\end{rebuttals}
\section{九、持方优劣评估}
\begin{assessment}
\assess{定义空间}{日常汉语里“用现实解释”天然带着“说实话”的味道，“用故事解释”则容易被听成“编一个说法哄他”。正方站在中性词上，反方要先花力气把“故事”从“哄骗”里摘出来才能开始论证，这在一辩稿里要占掉宝贵的二十秒}{偏正}
\assess{论证义务}{中立判准的两项指标里，认知准确这一项证据几乎一边倒。反方必须在这一项上先打平，才能靠情感这一项赢；正方只要在情感项上不输太多就能收下整场。反方的义务实质上更重}{偏正}
\assess{论据可得性}{发展心理学的实证研究（Slaughter 与 Lyons 的训练实验、Panagiotaki 的家庭对比、Ganea 的拟人化绘本实验）与所有权威机构的家长指南（美国儿科学会、NHS、Winston's Wish、Child Bereavement UK）都在正方一侧。反方拿得出的正面证据只有两条半：墨西哥的共存研究、叙事记忆的元分析，以及河南样本反映的“接收率不完全”。而且这两条都需要跨一步推理才能用上}{偏正}
\assess{评委直觉}{新生赛评委多半是学长学姐，很多人自己就是听着“姥姥变成星星了”长大的，对“跟四岁小孩讲心脏停止工作”有本能的不适。这是反方唯一的天然优势，而且不小}{偏反}
\assess{赛制顺序}{反四先结辩，可以写封路式结辩，把正四最可能走的那条路先点破，这对需要先扭转印象的反方有利；正四最后一发能真回应，对需要守住证据优势的正方也够用}{均衡偏反}
\end{assessment}
\begin{fields}
\claim[结论]{\textbf{反方较劣势}，劣势主要来自论据可得性与论证义务。正方手里有一个准实验、一个跨年龄段的家庭对比研究、一个绘本实验，外加四家权威机构的白纸黑字；反方手里最硬的一条是“父母的宗教与灵性思考对四个生物学子成分均无显著影响”，这是一个零效应，用来防守足够，用来进攻不够。}
\begin{sidecard}{劣势方打法}\end{sidecard}
\begin{enumerate}[leftmargin=1.8em,itemsep=2pt,topsep=2pt]
\item \textbf{收窄战场}。只打两件事：孩子几岁，解释的目的是什么。正方抛出的所有“故事害了孩子”的例子，一律用“那是敷衍不是故事”一句带回，不在细节上纠缠
\item \textbf{抢先定义}。反一立论的第一分钟必须把“故事”从“谎言”里摘出来，并且立刻给出《獾的礼物》这个锚：一个一句假话都没说的故事。这一步做成了，正方后面举的“睡着了”“弄丢了”就全部失效
\item \textbf{判准之争}。把“解释成功”的定义从“大人说了什么”改到“孩子接住了什么”。这是反方唯一能翻盘的支点，从一辩稿开始就要说，全场重复
\item \textbf{价值层}。反方的结辩是全场最有可能赢下印象分的一段。写一个具体的孩子，不要写“人类如何面对死亡”
\item \textbf{顺序利用}。反四先结辩，必须写成封路式：提前点破正四会说“我们不反对温柔，我们只反对不真实”，并在那句话落地之前先拆掉它
\end{enumerate}
\begin{sidecard}{优势方要防的}正方最容易输在被打成“冷酷派”。一旦评委心里形成“正方主张对四岁孩子讲器官衰竭”的画面，证据再多也救不回来。正方的防守动作有三个，缺一不可：一辩稿开场三十秒内说清“现实可以是温柔的，温柔不归故事独有”；主动承认故事作为辅助有价值，并且承认《獾的礼物》是好东西；全场每次引用数据之后补一个具体的孩子的画面，不要只报数字。\end{sidecard}
\end{fields}
\section{十、口径表}
\prosubsection{10.1 正方口径表}
\begin{canon}{}
\canongroup{定义}
\canonrow{现实}{以死亡的事实属性为解释的主要内容，核心是三句话：人死了不会再回来，身体停止运作，每个人最后都会死。语气可以温柔}{说清楚真正发生了什么；把事实讲在前面}{“冷冰冰地陈述医学事实”；“讲解剖学”}{本文档 3.1}
\canonrow{故事}{以叙事、比喻、拟人为主要载体的解释，字面内容不要求与事实一一对应}{用一个说法代替直说}{“故事就是骗孩子”；“童话就是谎言”}{本文档 3.2}
\canonrow{孩子}{三到十二岁，重心在三到八岁}{刚开始问死亡的孩子}{“小学生以上”；把三到五岁排除在外}{Speece 与 Brent 1984；Renaud 等 2015}
\canongroup{判准}
\canonrow{中立判准}{哪一种方式更能让孩子既正确地知道发生了什么，又承受得住它、并且愿意继续跟大人谈}{既要他懂，也要他受得住}{只说“哪种更准确”}{本文档第四节}
\canongroup{论点}
\canonrow{孩子问的是事实}{孩子来问死亡时问的是事实问题，用故事回答，他拿到的不是答案，是一个他会按字面理解的新说法}{答的是孩子没问的东西}{“所有故事都是答非所问”}{本文档 5.1}
\canonrow{先有现实才安心}{孩子是先有生物学理解，再用它约束比喻；而且越早有这份理解，越不怕死}{现实是地基，故事是盖在上面的楼}{“讲了现实孩子就不会怕”}{本文档 5.2}
\canongroup{数据}
\canonrow{提问错配}{墨西哥 61 名三岁半到六岁十一个月孩子里，关于死亡的一般提问 61.2\% 问的是死因，没有一个孩子问到宗教；家长的回答里 22.0\% 含宗教内容}{孩子问为什么会死，大人答去了天上}{把 22.0\% 说成“大多数家长”}{Gutiérrez 等 2020，《Child Development》，DOI 10.1111/cdev.13263}
\canonrow{训练实验}{60 名三岁七个月到五岁十一个月的孩子，40 人被教了身体如何维持生命，20 人对照；训练组对身体功能和对死亡的理解都显著提升}{教了身体，连带懂了死亡}{把准实验说成“随机对照试验”；说它“证明了该讲现实”}{Slaughter 与 Lyons 2003，《Cognitive Psychology》46(1): 1–30}
\canonrow{恐惧下降}{90 名四到八岁孩子，控制年龄与一般焦虑后，越成熟的生物学死亡理解对应越低的死亡恐惧}{懂得越清楚的孩子越不怕}{说成“讲现实降低了恐惧”（这是相关不是因果）}{Slaughter 与 Griffiths 2007，《Clinical Child Psychology and Psychiatry》12(4): 525–535}
\canonrow{家庭对比}{英国 136 人样本里，八到十一岁中宗教家庭孩子的生物学模型占 23\%，非宗教家庭占 77\%}{听来世长大的孩子，生物学理解落后一截}{把“父母信不信”直接说成“父母讲不讲故事”}{Panagiotaki 等 2018，《Journal of Experimental Child Psychology》166: 96–115}
\canonrow{拟人化绘本}{拟人化的动物绘本让三到五岁孩子学到的新事实更少，还更容易把人的特征安到真动物身上}{越可爱的书，教错的越多}{把它说成“关于死亡绘本的研究”}{Ganea 等 2014，《Frontiers in Psychology》5: 283}
\canonrow{长期影响}{33 位十八岁前失去父母的成年人的叙事研究：没有在合适时点按孩子的理解水平给出清楚诚实的信息，被认为对成年后的信任与自我价值有负面影响}{当年没人跟他说清楚，这件事跟了他一辈子}{说成“隐瞒会导致心理疾病”}{Ellis 等 2013，《Journal of the Royal Society of Medicine》106(2): 57–67}
\canongroup{案例}
\canonrow{睡着了}{美国儿科学会点名“睡着了”“去了一趟远门”这类说法，因为孩子按字面理解，可能因此害怕睡觉、害怕出门}{一个不敢关灯的孩子}{把机构指南说成“研究证明”}{HealthyChildren.org，Schonfeld 与 Nasir，2024 年 6 月 26 日更新}
\canonrow{一片叶子}{《一片叶子落下来》把死亡讲成“另一种形式的新生”，这句话和不可逆性正好相反}{那本很有名的绘本}{说这本书“有害”，改说“这句话对四岁孩子是误导”}{Buscaglia 1982}
\end{canon}
\consubsection{10.2 反方口径表}
\begin{canon}{}
\canongroup{定义}
\canonrow{故事}{以叙事、比喻、拟人为主要载体的解释，字面内容不要求与事实一一对应，但也不要求与事实冲突}{用一个孩子背得动的形状装住这件事}{“故事就是善意的谎言”；“故事可以不管真假”}{本文档 3.2}
\canonrow{现实}{以死亡的事实属性为解释的主要内容，语气可以温柔}{把生物学事实讲在前面}{“冷冰冰地陈述医学事实”；“正方要给孩子讲解剖”}{本文档 3.1}
\canonrow{孩子}{三到十二岁，重心在三到八岁；这个年纪的关键事实是死亡概念还在成形}{刚开始问死亡的孩子}{把讨论缩到只有三到五岁}{Speece 与 Brent 1984；Renaud 等 2015}
\canongroup{判准}
\canonrow{中立判准}{哪一种方式更能让孩子既正确地知道发生了什么，又承受得住它、并且愿意继续跟大人谈}{不看大人说了什么，看孩子接住了什么}{只说“哪种让孩子舒服”}{本文档第四节}
\canongroup{论点}
\canonrow{孩子先懂故事}{解释成功与否看孩子接住了什么；在他最需要解释的年纪，能被接住的形式是叙事}{讲得准，不等于讲得进去}{“孩子听不懂事实”（说成“接不住完整的生物学解释”）}{本文档 6.1}
\canonrow{故事多给一样}{孩子可以一边知道人死了不会回来，一边在故事里保留这段关系，而后一件事只有故事做得到}{一套用来知道，一套用来相处}{“故事比事实更真实”}{本文档 6.2}
\canongroup{数据}
\canonrow{成熟率}{河南 114 名五到六岁孩子，死亡概念五个子成分的成熟率是：不可逆性 67.8\%、因果性 60.7\%、非功能性 57.1\%、普遍性 55.3\%、死亡判定 53.6\%}{最容易的那一项也有三分之一的孩子还没掌握}{说成“中国孩子不懂死亡”}{Liu 与 Liu 2024，《Frontiers in Psychology》15，DOI 10.3389/fpsyg.2024.1376253}
\canonrow{共存}{墨西哥 61 名孩子里，81.9\% 相信亡者在亡灵节回来探望，78.6\% 相信亡者吃了供品；同时全部孩子都表现出至少部分生物学死亡理解}{既相信亡者回来，也知道死亡不可逆}{把它说成“故事提高了理解”}{Gutiérrez 等 2020}
\canonrow{零效应}{同一研究中，父母自报的宗教与灵性思考，对普遍性、终结性、非功能性、因果性四个子成分均无显著影响}{讲这些的家庭，孩子的生物学理解没被拖累}{把无显著影响说成“有正面作用”}{Gutiérrez 等 2020}
\canonrow{叙事元分析}{75 个以上独立样本、33000 多名参与者的元分析：故事比说明文更容易被理解，也被记得更牢}{同样一件事，讲成故事更进得去}{用它证明“故事讲的内容更对”；不说局限就用}{Mar 等 2021，《Psychonomic Bulletin \& Review》28(3): 732–749}
\canonrow{成人无共存}{Panagiotaki 等 2018 的成人组没有出现共存式思维，只有生物学解释与超自然解释的清楚区分}{听故事长大的人，后来分得很清楚}{说成“故事没有任何代价”}{Panagiotaki 等 2018}
\canongroup{案例}
\canonrow{獾的礼物}{獾死了，书里没让它复活，也没说它去了别处；朋友们靠回忆它教过的事活下去}{一本一句假话都没说的绘本}{把它说成“证明故事更准确”}{Varley 1984，《Badger's Parting Gifts》，Andersen Press}
\canonrow{清明}{一家人在墓前摆三只碗，说“爷爷你吃”；没有大人因此以为爷爷会回来，也没有孩子因此答错人会不会复活}{我们自己也在做同一件事}{说“清明证明了故事有效”}{中国民俗常识}
\canonrow{爷爷变成了幽灵}{爷爷死后变成幽灵回来，是因为忘了做一件事；说完再见，他就走了}{那本讲告别的绘本}{只讲前半段不讲结尾}{金·弗珀兹·艾克松文、爱娃·艾瑞克松图，2007 年中文版}
\end{canon}
\section{十一、论据出处清单}
\begin{sourcelist}
\source{1}{五到七岁多数儿童获得不可逆性、丧失功能性、普遍性的理解；四十余项研究的综述}{Speece M W, Brent S B (1984)，《Children's Concept of Death: A Review of Three Components of a Death Concept》，《Child Development》55(5): 1671–1686，ERIC EJ308874，\url{https://eric.ed.gov/?id=EJ308874}}{between five and seven years of age, the majority of healthy children in modern urban-industrial societies achieve an understanding of the irreversibility, nonfunctionality, and universality of death}{2026-09-18}{已核实}{双方定义、双方一辩}{「四十余项」的篇数取自 Liu 与 Liu 2024 对该综述的描述（by analyzing more than 40 studies），1984 年原文正文本次未打开，只核到 ERIC 著录的摘要结论}
\source{2}{60 名三岁七个月至五岁十一个月儿童，40 人受训、20 人对照；训练提升身体功能与死亡理解}{Slaughter V, Lyons M (2003)，《Learning about life and death in early childhood》，《Cognitive Psychology》46(1): 1–30，DOI 10.1016/s0010-0285(02)00504-2，PubMed 12646154，\url{https://pubmed.ncbi.nlm.nih.gov/12646154/}}{Sixty preschool children between the ages of 3 years, 7 months and 5 years, 11 months participated……The Control group (n=20) received no training}{2026-09-18}{已核实}{正方论点二}{因果结论对应同一摘要的 this learning coincided with significant development in their understanding of human body function, and of death}
\source{3}{90 名四到八岁儿童；控制年龄与一般焦虑后，生物学死亡理解越成熟，死亡恐惧越低}{Slaughter V, Griffiths M (2007)，《Death understanding and fear of death in young children》，《Clinical Child Psychology and Psychiatry》12(4): 525–535，DOI 10.1177/1359104507080980，PubMed 18095535，\url{https://pubmed.ncbi.nlm.nih.gov/18095535/}}{A regression analysis indicated that more mature death understanding was associated with lower levels of death fear, when age and general anxiety were controlled}{2026-09-18}{已核实}{正方论点二}{}
\source{4}{死亡概念四成分；五岁前普遍性、五岁身体机能停止、六岁因果性；儿童先形成生物学理解再约束宗教理解}{Menendez D, Hernandez I G, Rosengren K S (2020)，《Children's Emerging Understanding of Death》，《Child Development Perspectives》，DOI 10.1111/cdep.12357，作者自存全文 \url{https://rosengrenlab.digitalscholar.rochester.edu/wp-content/uploads/2020/06/MenendezHernandezRosengren2020.pdf}}{By age 5, most children understand that death involves the cessation of bodily processes, and by age 6, children have the more sophisticated understanding that death can be caused by many factors, not just old age}{2026-09-18}{已核实}{双方定义、正方论点二}{「先生物学再约束宗教」对应同文 children come to a biological understanding of death prior to integrating spiritual or religious dimensions}
\source{5}{墨西哥 53 家庭 61 名 3.5–6.9 岁儿童；提问 61.2\% 问死因、无人问宗教，家长回答 22.0\% 含宗教；81.9\% 信亡者探望、78.6\% 信亡者进食；全部孩子有部分生物学理解；父母宗教与灵性思考对四子成分均无显著影响；31.1\% 的孩子在亡灵节感到害怕（4 岁 55.0\%、5 岁 28.6\%、6 岁 15.0\%）}{Gutiérrez I T 等 (2020)，《Embracing Death: Mexican Parent and Child Perspectives on Death》，《Child Development》91(2): e491–e511，DOI 10.1111/cdev.13263，作者自存全文 \url{https://rosengrenlab.digitalscholar.rochester.edu/wp-content/uploads/2020/06/Gutierrez_et_al-2019-Child_Development.pdf}}{Fifty children (81.9\%) stated that they believed their dead relatives came back to visit……48 (78.6\%) believed that the dead relatives consumed the food}{2026-09-18}{已核实}{正方论点一、反方论点二}{同文其余数字的位置：61.2\%／22.0\% 见 Discussion 关于提问与回答主题的段落；31.1\%（4 岁 55.0\%、5 岁 28.6\%、6 岁 15.0\%）见 Results 亡灵节情绪反应段；零效应见 Results 各子成分段（普遍性 p = .799、终结性 p = .751、非功能性 p = .698、因果性 p = .365）；全部孩子有生物学理解见摘要 while all children in this sample displayed a biological understanding of death}
\source{6}{英国 136 人；6–7 岁与 8–9 岁非宗教家庭孩子在普遍性上更常生物学正确（1.67 比 0.83；1.65 比 0.89）；8–11 岁宗教家庭生物学模型 23\% 对 77\%；成人组无共存式思维；受访成人承认自己不信来世却对孩子讲“奶奶在天上还看得见你”}{Panagiotaki G, Hopkins M, Nobes G, Ward E, Griffiths D (2018)，《Children's and adults' understanding of death: Cognitive, parental, and experiential influences》，《Journal of Experimental Child Psychology》166: 96–115，DOI 10.1016/j.jecp.2017.07.014；核实用的是 UEA 机构库作者接受稿全文 \url{https://ueaeprints.uea.ac.uk/64305/1/Accepted_manuscript.pdf}}{6-7 year-olds of nonreligious parents were more frequently biologically correct than those of religious parents (M=1.67 vs M=.83)}{2026-09-18}{已核实}{正方论点二、反方论点二}{同一作者稿其余位置：8–9 岁 M=1.65 vs M=.89 紧接上句；23\% vs 77\% 见 fewer biological (23\% vs. 77\%) models 一句；成人无共存见摘要 There was no evidence for coexistent thinking among adults；家长原话见 Discussion 的 your granny can still see you now that she is in the sky。核实依据是作者接受稿，未打开 Elsevier 正式排版版}
\source{7}{拟人化绘本降低三到五岁儿童对新事实的学习，并增加对真实动物的拟人化归因}{Ganea P A, Canfield C F, Simons-Ghafari K, Chou T (2014)，《Do cavies talk? The effect of anthropomorphic picture books on children's knowledge about animals》，《Frontiers in Psychology》5: 283，DOI 10.3389/fpsyg.2014.00283，PMC3989584，\url{https://pmc.ncbi.nlm.nih.gov/articles/PMC3989584/}}{These results indicate that anthropomorphized animals in books may not only lead to less learning but also influence children's conceptual knowledge of animals}{2026-09-18}{已核实}{正方反驳库 3}{}
\source{8}{75 个以上独立样本、33000 多名参与者的元分析：故事比说明文更易理解、记得更牢}{Mar R A, Li J, Nguyen A T P, Ta C P (2021)，《Memory and comprehension of narrative versus expository texts: A meta-analysis》，《Psychonomic Bulletin \& Review》28(3): 732–749，DOI 10.3758/s13423-020-01853-1，PMC8219577，\url{https://pmc.ncbi.nlm.nih.gov/articles/PMC8219577/}}{Based on over 75 unique samples and data from more than 33,000 participants, we found that stories were more easily understood and better recalled than essays}{2026-09-18}{已核实}{反方论点一}{}
\source{9}{33 位十八岁前丧亲成年人的叙事研究：未在合适时点给出清楚诚实的信息，对成年后信任与自我价值有负面影响}{Ellis J, Dowrick C, Lloyd-Williams M (2013)，《The long-term impact of early parental death: lessons from a narrative study》，《Journal of the Royal Society of Medicine》106(2): 57–67，DOI 10.1177/0141076812472623，PMC 免费全文 \url{https://pmc.ncbi.nlm.nih.gov/articles/PMC3569022/}}{failure to provide clear and honest information at appropriate time points relevant to the child's level of understanding was perceived to have a negative impact in adulthood}{2026-09-18}{已核实}{正方论点二、正方反驳库 5}{样本描述见同文 33 adults (7 men and 26 women) who had experienced parental death during childhood}
\source{10}{避免“睡着了”“去了很远的地方”，因为孩子按字面理解；五岁以下不一定完全理解，五到七岁一般学会这些概念}{美国儿科学会 HealthyChildren.org，《How Children Understand Death: What to Say When a Loved One Dies》，Schonfeld D J 与 Nasir A 撰，2024 年 6 月 26 日更新，\url{https://www.healthychildren.org/English/healthy-living/emotional-wellness/Building-Resilience/Pages/How-Children-Understand-Death-What-You-Should-Say.aspx}}{avoid using vague phrases such as, “Your grandparent has gone to sleep (or gone on a long journey).”}{2026-09-18}{已核实}{正方论点一、反方论点一}{年龄口径对应同页 Children under 5 years old may not grasp the full message 与 children generally learn this concepts between 5-7 years of age on average}
\source{11}{直接、诚实、开放的沟通比隐瞒真相更有帮助}{英国国民保健署 NHS，《Children and bereavement》，最后审阅 2023 年 1 月 9 日，\url{https://www.nhs.uk/mental-health/children-and-young-adults/advice-for-parents/children-and-bereavement/}}{Direct, honest and open communication is more helpful than trying to protect your child by hiding the truth}{2026-09-18}{已核实}{正方反驳库 7}{}
\source{12}{河南 114 名五到六岁儿童五子成分成熟率：67.8\% / 60.7\% / 57.1\% / 55.3\% / 53.6\%；86\% 的家庭会公开谈论死亡}{Liu J, Liu F (2024)，《Research on the development of the concept of death in children and its influencing factors》，《Frontiers in Psychology》15，DOI 10.3389/fpsyg.2024.1376253，\url{https://www.frontiersin.org/journals/psychology/articles/10.3389/fpsyg.2024.1376253/full}}{children's “universality of death” fully mature accounted for 55.3\%……60.7\% of children's “Causality of death” has reached full maturity}{2026-09-18}{已核实}{反方论点一、反方反驳库 7}{同文其余数字的位置：不可逆性 67.8\% 见 Table 9，非功能性 57.1\% 见 Table 10（原文作 Inorganicity），死亡判定 53.6\% 见 Table 13；86\% 见 Table 15 与 Ninety-eight (98) or 86\% had talked openly about the topic of death 一句}
\source{13}{报道本身：中国死亡教育的缺位现状}{澎湃新闻·有数，《这堂人生必修课，大多数中国人都还欠缺》，曹若昕、马昊宁、逄博、吴佳卉、杨柠瑜，2024 年 4 月 4 日，\url{https://www.thepaper.cn/newsDetail_forward_26912005}}{死亡教育还没有找到属于自己的位置，中小学的死亡教育内容往往渗透在心理健康教育、思想道德教育中}{2026-09-18}{已核实}{背景，不上场}{}
\source{14}{精细化回忆叙事风格与儿童自传体记忆、情绪理解、自我概念的关系}{Fivush R, Haden C A, Reese E (2006)，《Child Development》77(6): 1568–1588（书目已核实，内容未开原文）}{}{}{有把握}{反方论点一备用学理}{查 DOI 10.1111/j.1467-8624.2006.00960.x 原文}
\source{15}{两种思维模式：范式模式与叙事模式，互不可还原}{Bruner J (1986)，《Actual Minds, Possible Worlds》，Harvard University Press}{}{}{有把握}{反方论点一学理}{查该书第二章 Two Modes of Thought}
\source{16}{持续联结：与逝者维持联结不是病理，而是正常且具保护性的哀伤方式}{Klass D, Silverman P R, Nickman S (1996)，《Continuing Bonds: New Understandings of Grief》}{}{}{有把握}{反方论点二学理}{查原书导论与 Silverman 关于丧亲儿童的章节}
\source{17}{用“死了”而不用“走了”“弄丢了”；孩子会问“如果妈妈丢了，为什么我们不去找她”}{Winston's Wish，《Telling a child someone has died》；Child Bereavement UK，《Explaining to a child that someone has died》（页面本次访问返回 403，依据检索摘要）}{}{}{有把握}{正方论点一案例}{换网络环境直接打开两家官网页面抄原句}
\source{18}{加拿大网络调查：75\% 的家长谈过死亡，最早的对话发生在孩子三到三岁半}{Renaud S J, Engarhos P, Schleifer M, Talwar V (2015)；经 Gutiérrez 等 2020 原文转引}{}{}{有把握}{双方年龄口径}{查原文《Children's Earliest Experiences with Death》}
\source{19}{83\% 的家长说孩子问过死亡；儿童最爱读物与凯迪克奖绘本中只有 3\% 描写死亡}{Rosengren K S 等 (2014)；经 Menendez 等 2020 与 Gutiérrez 等 2020 转引}{}{}{有把握}{背景、正方反驳库 8}{查 Rosengren 等 2014 专著原文}
\source{20}{75\% 的儿童动画电影里有死亡情节}{Bridgewater E, Menendez D, Rosengren K S (2019)，未发表手稿；经 Menendez 等 2020 转引}{}{}{有把握}{正方淘汰论点}{该手稿可能已正式发表，查 Menendez 后续论文}
\source{21}{回忆父母当年开放谈论死亡的人，童年丧亲后应对更好，并延续到成年}{Martincekova L 等 (2018)；经 Menendez 等 2020 转引}{}{}{有把握}{正方备用数据}{查原文期刊与样本}
\source{22}{绘本书目：《一片叶子落下来》(Buscaglia, 1982)、《獾的礼物》(Varley, 1984, 1985 年 Mother Goose Award)、《爷爷变成了幽灵》(艾克松文、艾瑞克松图, 2007 年中文版)}{各书豆瓣与出版社页面、Andersen Press 官网}{}{}{有把握}{双方案例}{核对中译本出版社与年份}
\source{23}{库伯勒·罗斯五阶段模型缺乏实证支持，且不适用于儿童}{多篇综述与科普批评文章}{}{}{有把握}{双方反驳库备用}{查 Stroebe 等对阶段模型的实证检验论文}
\source{24}{1031 人调查：46.43\% 自认缺乏面对死亡的知识，85.07\% 从未主动进行死亡咨询}{澎湃新闻转引北医三院宋子妹等的研究，原始论文未找到}{}{}{待核实}{不要上场使用}{在知网搜“宋子妹 死亡教育 调查”}
\source{25}{教育部 2021 年 11 月《生命安全与健康教育进中小学课程教材指南》，5 个领域 30 个核心要点}{教育部政府门户网站（本次访问遇重定向，未打开）}{}{}{待核实}{中国政策背景}{直接访问 moe.gov.cn 该文号页面}
\source{26}{部编版初中《道德与法治》涉及死亡的内容不到两页}{网络转述，无原始出处}{}{}{待核实}{不要上场使用}{直接翻教材目录逐页核对}
\source{27}{289 位中国父母“死亡话语素养”调查}{《与儿童谈生论死》，华人生死学期刊（本次访问返回 503）}{}{}{待核实}{中国家长现状}{换时间重试该期刊网站或查知网}
\end{sourcelist}
\textbf{上场纪律}：状态为【待核实】的四条，赛前没查清就\textbf{不准出现在任何一篇稿子里}。标“有把握”的可以用，但引用时只报机构与年份，不报具体数字；被追问时如实说“这是我方从综述中读到的转引，原始论文我方赛前未能取得”。

\section{十二、备赛建议}
\begin{fields}
\begin{timeline}[时间表]
\when{距比赛 7 天}{全队读完本文档第三、四节，两位一辩背下中立定义与中立判准的原话；分工把第十一节的四条【待核实】查完}
\when{距比赛 5 天}{两位二辩、三辩、四辩各自对着第七、八节的反驳库做一遍“我怎么回”的口头演练；二辩与四辩把第七节表格里的每一条答案对齐}
\when{距比赛 3 天}{三份文稿全部出声读一遍并计时，超时按“论据 → 修饰语 → 机制中间步骤”删}
\when{距比赛 1 天}{完整模辩一场，重点看下面三件事}
\end{timeline}
\field{模辩重点检验}{
\begin{enumerate}[leftmargin=1.8em,itemsep=2pt,topsep=2pt]
\item 正方有没有在开场三十秒内摆脱“冷酷派”的印象；反方有没有在开场三十秒内把“故事”从“谎言”里摘出来
\item 三辩盘问对方二辩与四辩同一个问题时，两人的答案是不是同一个意思。最可能露馅的两个问题是“你们说的孩子几岁”和“亲人今天刚去世也这样吗”
\item 那一次硬碰硬（Panagiotaki 的家庭对比对上 Gutiérrez 的零效应）双方是不是都能在二十秒内说清楚“对方那份研究为什么推不倒我”
\end{enumerate}
}
\begin{checklist}[每位队员赛前要背下来的]
\item 我方的中立定义原话、中立判准原话、两个论点标签、对方两个论点各一句话反驳、我方那条最强数据的机构与年份、以及一辩稿里抛出的那个贯穿全场的问题。
\end{checklist}
\end{fields}
\end{document}
